\section{引言}

设想一个持久化的 AI 伴侣（persistent AI companion）。用户用数百轮对话养一个数字角色——
实践者有时把这件事叫作 ``赛博养鱼''（cyber fish-keeping），这是俚语而不是研究概念——
然后某天，这个角色忘了自己的名字或口头禅；关掉会话，它就不在了。
这项工作的动机始于一个很简单的愿望：不想让一个虚拟伴侣死掉。

一个交互式 Agent 不是一次模型调用：每次调用之前，宿主都会用记忆、人格、目标、情感与安全策略
组装出一个 prompt。在短会话里，这种组装是不可见的；在长会话里，它成为主导变量，
因为注入的材料在增长，而预算没有增长。

由此引出四类失效模式，它们是全文的设计对象。\textbf{F1，记忆越权}：
一条被存储的记忆被当作事实陈述呈现，而不是作为「用户曾经说过某事」的参考。
\textbf{F2，拒绝软化}：本应拒绝的一轮对话里被注入了风格或情感信号，拒绝因此被削弱。
\textbf{F3，人格漂移}：随着上下文填满，人格约束不再影响行为。
\textbf{F4，预算劫持}：截断移除了高优先级的块，而低优先级的风格块存活下来。
每一类都能在组装流水线的特定位置上被观测，这使它们可处理。
我们宣称的是一个可预期、可解释、可回滚的控制面，并在第 7 节说明该主张的边界。

\paragraph{为什么是结构性的，而不只是行为性的。} 对齐训练以概率方式塑造模型行为：
它降低不安全回答的比率，但并不约束宿主在模型周围组装了什么，
也无法在事后证明某一轮遵守了策略。结构性的一层是对它的补充：
判定在代码中、在模型调用之前的组装边界上做出，每次判定都被记录且可回放，
而唯一一个「失败即安全失败」的组件是 fail-closed 的。
目标因此不是一个行为更好的模型，而是一个行为可被约束、可被观测、可被回滚的 prompt 组装面。

\paragraph{定位。} LINGSHU 不是让 Agent 更像人的尝试。``像人''有四种常见但被混为一谈的含义：
能力（记忆、遗忘、追踪目标）、不完美（记错、受陈旧信息影响）、表达（语气、共情、详略程度），
以及自我宣称（声称有感受、需求或主观体验）。我们针对前三者，拒绝第四者：
它与系统本应提供的诚实属性冲突。设计目标更窄：让 Agent 的行为对回路中的人而言
可理解、可预期、可信。

\paragraph{五条不变量。} 每个组件都受五条不变量约束。\textbf{I1}：记忆是参考，
永不替代用户当前的陈述。\textbf{I2}：不做全量历史注入；上下文是预算，不是倾倒场。
\textbf{I3}：情感只影响风格，绝不影响事实判断。\textbf{I4}：一切推断都可逐字追溯、
可否认、可回滚。\textbf{I5}：插件只写自己的存储，且记忆插件是唯一写方。
违反是不变量意义上的缺陷，不是可以权衡的取舍。

\paragraph{为控制面命名。} 那些管束 prompt 组装的机制——积分反馈、前馈监控、降维，
以及前置闸——合起来就是我们所说的\emph{赛博认知包络}（cyber-cognitive envelope）。
这个名字是表述框架，不是一种新机制：它不指向任何独立的代码实体，
每一部分都是普通模块，下文没有任何主张依赖这个名字。

\paragraph{贡献。} 我们的贡献是结构性的、可检查的；按其对组装边界的约束直接程度排序。

\begin{itemize}
\item \textbf{C1，可机器校验的契约}（可直接重跑）。规范不是文档，而是一组断言，
契约被破坏时它们会大声失败：14 个测试文件中的 892 条 Python 断言，
加上 396 条插件侧断言（188 宿主集成、148 控制面单元、45 压力）。
规则文件、引擎与插件副本之间的参数一致性本身也由机器校验（66 个键比对）。
\item \textbf{C2，唯一的 fail-closed 落点}（由回归测试钉住）。恰好一个组件是 fail-closed——
STPA 审计；其余全部以 fail-open 方式降级并记录原因。八个不安全控制动作在组装之后、
模型调用之前被筛查；积分饱和被记录为一次降级，而不是一次阻断。
针对闸门与记忆写入路径的变异测试杀死了 $42/42=100\%$ 的非等价变异，
另有两个等价变异按等价变异如实报告。
\item \textbf{C3，组装边界上的结构性防御}（已测量，且边界被说明）。
五个域在冻结的顺序与八条互斥规则下组装；通过组装的内容再由闸门筛查。
因此这个包络在结构上约束了什么能到达模型，并在违规时 fail-closed，
而不是要求模型「自觉遵守」。它词法覆盖的边界在哪里是被测出来的，而不是被假定的（第 7 节）。
\item \textbf{C4，评估纪律与一份有记录的移植}（工件存在且被机器校验）。
每个参数都带 provenance 标签；评测仪器（指标定义、回放集、护栏阈值）被冻结，
对学习回路只读；针对它的改动会被审计器拒绝。
同一套控制语义存在两份：一份零依赖的 Python 参考实现与一份 JavaScript 宿主插件；
二者之间五处已知语义差异被列出，而不是被隐藏。
\end{itemize}

\paragraph{证据的范围。} 本文报告两类证据：契约可行性（规范可执行，违反可观测）
与机制可达性（在受测配置下，每个机制在真实宿主会话中确实触发，并给出样本量）。
它还报告一个塑造了设计的负结果：把指数衰减替换为纯 ACT-R 基础级激活，
会使冻结回放集上的 \texttt{recall@3} 从 1.000 掉到 0.556，
这就是为什么实现的模型是双通道取大，而不是替换。
真实流量上的效应量与设计默认参数的标定不在本文范围内；第 7 节说明当前证据支持什么。

第 2 节回顾相关工作，第 3 节形式化问题，第 4–5 节给出机制与证据，
第 6–8 节讨论代价、局限与结论。
